
\documentclass[12pt,a4paper]{article}

\usepackage[a4paper,margin=1in]{geometry}
\usepackage{graphicx}
\usepackage{booktabs}
\usepackage{amsmath,amssymb}
\usepackage{float}
\usepackage{hyperref}
\usepackage{xcolor}
\usepackage{listings}
\usepackage{enumitem}
\usepackage{fancyhdr}
\usepackage{longtable}

\hypersetup{colorlinks=true,linkcolor=blue,urlcolor=blue}

\pagestyle{fancy}
\fancyhf{}
\lhead{ICS1512}
\rhead{Machine Learning Algorithms Laboratory}
\cfoot{\thepage}

\lstset{
language=Python,
basicstyle=\ttfamily\small,
numbers=left,
frame=single,
breaklines=true,
showstringspaces=false
}

\begin{document}

\begin{tabular}{|p{4cm}|p{11cm}|}
\hline
Experiment No. & 7 \\ \hline
Experiment Title & Dimensionality Reduction and Model Evaluation (With and Without PCA)\footnotemark \\ \hline
Student Name & Sanjay S \\ \hline
Register Number & 3122247001056 \\ \hline
Faculty & Dr. Poreddy Ajay Kumar Reddy \\ \hline
Submission Date & 21/09/2026 \\ \hline
\end{tabular}
\footnotetext{\textbf{GitHub Repository: }\url{https://github.com/sanjay-6727/Introdution-to-Machine-Learning-}}

\section{Objective}
See what PCA actually does to model performance instead of just assuming it helps. The plan was to
take the same ten classifiers, run each one on the raw 10-feature space and again on a PCA-reduced
space, tune both properly, and compare like for like using 5-fold CV so any difference in the numbers
is actually about PCA and not about one side getting a better hyperparameter search.

\section{Problem Statement}
Given a three-class classification dataset, standardize everything, fit ten different classifiers
under two conditions, No-PCA and With-PCA, tuning each with a grid search and validating with 5-fold
cross-validation plus a held-out test set. Then compare accuracy, F1 and how stable each model is
across folds, and figure out which models actually benefit from dimensionality reduction and which
ones just get worse.

\section{Dataset Description}
\begin{longtable}{|p{3cm}|p{3.6cm}|p{1.6cm}|p{1.8cm}|p{2cm}|p{2.5cm}|}
\hline
\textbf{Dataset} & \textbf{Source} & \textbf{Samples} & \textbf{Features} & \textbf{Target Classes} & \textbf{Missing Values} \\ \hline
PCA Classification Dataset & Provided course dataset, \texttt{pca\_classification\_clean\_1000.csv} & 1000 & 10 numeric features (\texttt{feature\_1}...\texttt{feature\_10}) & 3 (Class\_0, Class\_1, Class\_2) & none \\ \hline
\end{longtable}

\textbf{Class distribution:} Class\_0 = 333, Class\_1 = 333, Class\_2 = 334, basically a perfect
three-way split, so no resampling or class-weighting was needed anywhere. The file also has a
continuous \texttt{target} column sitting alongside \texttt{target\_class}, which I dropped since the
job here is classification, not regression.

\textbf{Preprocessing steps:}
\begin{itemize}
\item Dropped the raw \texttt{target} column, keeping \texttt{target\_class} as the label.
\item Label-encoded \texttt{target\_class} into integers 0, 1, 2 with \texttt{LabelEncoder}.
\item No missing values were present, so no imputation was required.
\item An 80/20 train-test split was used, stratified on the class label (800 train, 200 test).
\item All 10 features were standardized with \texttt{StandardScaler}, fit on the training split only,
inside the pipeline so it is refit correctly within every cross-validation fold.
\end{itemize}

\section{PCA Design and Scree Plot}
Rather than picking one fixed number of components up front, I let each model's grid search try
component counts from 2 through 8 and pick whatever worked best for it, since there's no real reason
every classifier should want the same amount of compression. Once I plotted the actual variance curve
on the standardized training data, though, a pattern shows up:

\begin{table}[H]
\centering
\caption{PCA Variance Explained}
\begin{tabular}{p{3.5cm}p{4.5cm}p{2.7cm}p{4.5cm}}
\toprule
Setting & Chosen Components / Variance Target & Explained Variance (\%) & Justification \\
\midrule
With-PCA (per-model tuned) & 2--8 components searched per model; most models settled on 3 & 86.77\% (3 comps) & Grid search kept landing on 3 components pretty much across the board, and it's not hard to see why -- the first three PCs each carry roughly 27--30\% of the variance on their own, so together they already cover most of what matters. Technically you need 4 components to clear the usual 95\% cutoff, but since 3 was what actually won on accuracy for almost every model, I went with that as the reference point instead of chasing the textbook threshold. \\
\bottomrule
\end{tabular}
\end{table}

\begin{figure}[H]
\centering
\includegraphics[width=0.75\linewidth]{outputs/scree_plot.png}
\caption{Scree plot: individual and cumulative explained variance per principal component. The first
3 components capture about 86.8\% of the total variance, and 4 components are needed to cross the
95\% threshold.}
\end{figure}

\section{Machine Learning Algorithms and Hyperparameter Grids}
Ten classifiers went through this: SVM, Naive Bayes, KNN, Logistic Regression, Decision Tree, Random
Forest, AdaBoost, Gradient Boosting, XGBoost, and a Stacking ensemble on top of all that. Each one sat
inside a \texttt{Pipeline} -- \texttt{StandardScaler} then an optional \texttt{PCA} step then the
actual estimator -- and \texttt{GridSearchCV} handled the tuning with 5-fold CV and accuracy as the
scoring metric, so the whole thing stayed consistent from model to model.

\begin{table}[H]
\centering
\caption{SVM -- Hyperparameter Tuning Results}
\begin{tabular}{p{2.2cm}p{2.6cm}p{2.6cm}p{2.6cm}p{2.6cm}}
\toprule
Kernel Values Tried & C Values Tried & Gamma Values Tried & Performance (No-PCA) & Performance (With-PCA) \\
\midrule
linear, rbf & 0.1, 1, 10 & scale, auto & Best: linear, C=0.1, gamma=scale; Test Acc=0.860, F1=0.8616 & Best: linear, C=10, gamma=scale, PCA=3; Test Acc=0.835, F1=0.8356 \\
\bottomrule
\end{tabular}
\end{table}

\begin{table}[H]
\centering
\caption{Naive Bayes -- Smoothing Choices}
\begin{tabular}{p{4.5cm}p{4.5cm}p{4.5cm}}
\toprule
Smoothing Parameter (var\_smoothing) & Performance (No-PCA) & Performance (With-PCA) \\
\midrule
$1\times10^{-11}, 10^{-10}, 10^{-9}, 10^{-8}, 10^{-7}$ & Best: $1\times10^{-11}$; Test Acc=0.770, F1=0.7738 & Best: $1\times10^{-11}$, PCA=7; Test Acc=0.720, F1=0.7270 \\
\bottomrule
\end{tabular}
\end{table}

\begin{table}[H]
\centering
\caption{KNN -- Hyperparameter Tuning}
\begin{tabular}{p{2.3cm}p{2.3cm}p{2.6cm}p{3cm}p{3cm}}
\toprule
k Values Tried & Weights & Distance Metrics & Performance (No-PCA) & Performance (With-PCA) \\
\midrule
3, 5, 7, 9, 11 & uniform, distance & euclidean, manhattan & Best: k=9, distance, euclidean; Test Acc=0.820, F1=0.8201 & Best: k=11, distance, euclidean, PCA=3; Test Acc=0.845, F1=0.8455 \\
\bottomrule
\end{tabular}
\end{table}

\begin{table}[H]
\centering
\caption{Remaining Models -- Hyperparameter Grids and Best Settings}
\begin{longtable}{p{3cm}p{5.3cm}p{3.1cm}p{3.1cm}}
\toprule
Model & Grid Searched & Best (No-PCA) & Best (With-PCA) \\
\midrule
Logistic Regression & C: 0.01, 0.1, 1, 10, 100 & C=0.1; Acc=0.850, F1=0.8514 & C=0.01, PCA=3; Acc=0.840, F1=0.8395 \\
Decision Tree & max\_depth: None,3,5,10,20; min\_samples\_split: 2,5,10; criterion: gini, entropy & entropy, depth=None, split=5; Acc=0.795, F1=0.7957 & gini, depth=5, split=10, PCA=3; Acc=0.790, F1=0.7909 \\
Random Forest & n\_estimators: 100,200; max\_depth: None,5,10; min\_samples\_split: 2,5 & depth=None, split=2, n=200; Acc=0.815, F1=0.8143 & depth=10, split=5, n=100, PCA=3; Acc=0.825, F1=0.8250 \\
AdaBoost & n\_estimators: 50,100,200; learning\_rate: 0.01,0.1,1.0 & lr=1.0, n=200; Acc=0.795, F1=0.8005 & lr=1.0, n=200, PCA=5; Acc=0.800, F1=0.8041 \\
Gradient Boosting & n\_estimators: 50,100; learning\_rate: 0.05,0.1; max\_depth: 2,3,5 & lr=0.1, depth=3, n=100; Acc=0.785, F1=0.7847 & lr=0.05, depth=3, n=100, PCA=3; Acc=0.825, F1=0.8244 \\
XGBoost & n\_estimators: 50,100; max\_depth: 3,5; learning\_rate: 0.05,0.1 & lr=0.1, depth=5, n=100; Acc=0.790, F1=0.7903 & lr=0.1, depth=3, n=50, PCA=3; Acc=0.800, F1=0.7972 \\
Stacking & Base: linear SVM (C=10), Random Forest (n=100, depth=10), XGBoost (n=100, depth=3); Meta: Logistic Regression, 5-fold internal blending & Acc=0.855, F1=0.8565 & PCA=3; Acc=0.840, F1=0.8410 \\
\bottomrule
\end{longtable}
\end{table}

\section{Cross-Validation Results (All Models)}
\begin{table}[H]
\centering
\caption{5-Fold Cross-Validation Results -- No-PCA}
\begin{tabular}{lccccccc}
\toprule
Model & Fold 1 & Fold 2 & Fold 3 & Fold 4 & Fold 5 & Avg & Std \\
\midrule
SVM & 0.8938 & 0.8813 & 0.8750 & 0.8625 & 0.8313 & 0.8688 & 0.0213 \\
Naive Bayes & 0.8188 & 0.8625 & 0.8063 & 0.8375 & 0.7688 & 0.8188 & 0.0314 \\
KNN & 0.8500 & 0.8938 & 0.8125 & 0.8500 & 0.8313 & 0.8475 & 0.0270 \\
Logistic Regression & 0.8875 & 0.8938 & 0.8750 & 0.8625 & 0.8375 & 0.8712 & 0.0200 \\
Decision Tree & 0.7875 & 0.7813 & 0.7688 & 0.8063 & 0.7500 & 0.7788 & 0.0188 \\
Random Forest & 0.8750 & 0.9000 & 0.8438 & 0.8625 & 0.8313 & 0.8625 & 0.0240 \\
AdaBoost & 0.7750 & 0.8375 & 0.7813 & 0.8313 & 0.7625 & 0.7975 & 0.0308 \\
Gradient Boosting & 0.8063 & 0.8688 & 0.7938 & 0.8250 & 0.8125 & 0.8212 & 0.0258 \\
XGBoost & 0.8375 & 0.8313 & 0.8438 & 0.8313 & 0.8125 & 0.8312 & 0.0105 \\
Stacking & 0.9063 & 0.8688 & 0.8563 & 0.8313 & 0.8875 & 0.8700 & 0.0257 \\
\bottomrule
\end{tabular}
\end{table}

\begin{table}[H]
\centering
\caption{5-Fold Cross-Validation Results -- With-PCA}
\begin{tabular}{lccccccc}
\toprule
Model & Fold 1 & Fold 2 & Fold 3 & Fold 4 & Fold 5 & Avg & Std \\
\midrule
SVM & 0.8938 & 0.9000 & 0.8813 & 0.8688 & 0.8438 & 0.8775 & 0.0200 \\
Naive Bayes & 0.8438 & 0.8625 & 0.8563 & 0.8250 & 0.8125 & 0.8400 & 0.0188 \\
KNN & 0.9063 & 0.9125 & 0.8188 & 0.8500 & 0.8375 & 0.8650 & 0.0376 \\
Logistic Regression & 0.9063 & 0.8875 & 0.8688 & 0.8938 & 0.8188 & 0.8750 & 0.0306 \\
Decision Tree & 0.8188 & 0.8500 & 0.8563 & 0.8688 & 0.8063 & 0.8400 & 0.0236 \\
Random Forest & 0.8500 & 0.9125 & 0.8625 & 0.8875 & 0.8563 & 0.8738 & 0.0232 \\
AdaBoost & 0.8313 & 0.8563 & 0.8063 & 0.8313 & 0.8313 & 0.8312 & 0.0158 \\
Gradient Boosting & 0.8313 & 0.9063 & 0.8563 & 0.8875 & 0.8625 & 0.8688 & 0.0259 \\
XGBoost & 0.8438 & 0.9063 & 0.8688 & 0.8938 & 0.8188 & 0.8662 & 0.0320 \\
Stacking & 0.9000 & 0.8563 & 0.8813 & 0.8375 & 0.8938 & 0.8738 & 0.0235 \\
\bottomrule
\end{tabular}
\end{table}

\section{Final Comparison and ANOVA}
\begin{table}[H]
\centering
\caption{Final PCA vs No-PCA Comparison}
\resizebox{\textwidth}{!}{
\begin{tabular}{lcccccccc}
\toprule
Model & No-PCA CV & PCA CV & No-PCA Acc & PCA Acc & No-PCA F1 & PCA F1 & Acc Change & F1 Change \\
\midrule
SVM & 0.8688 & 0.8775 & 0.860 & 0.835 & 0.8616 & 0.8356 & -0.025 & -0.0260 \\
Naive Bayes & 0.8188 & 0.8400 & 0.770 & 0.720 & 0.7738 & 0.7270 & -0.050 & -0.0468 \\
KNN & 0.8475 & 0.8650 & 0.820 & 0.845 & 0.8201 & 0.8455 & +0.025 & +0.0254 \\
Logistic Regression & 0.8712 & 0.8750 & 0.850 & 0.840 & 0.8514 & 0.8395 & -0.010 & -0.0119 \\
Decision Tree & 0.7788 & 0.8400 & 0.795 & 0.790 & 0.7957 & 0.7909 & -0.005 & -0.0049 \\
Random Forest & 0.8625 & 0.8738 & 0.815 & 0.825 & 0.8143 & 0.8250 & +0.010 & +0.0107 \\
AdaBoost & 0.7975 & 0.8312 & 0.795 & 0.800 & 0.8005 & 0.8041 & +0.005 & +0.0035 \\
Gradient Boosting & 0.8212 & 0.8688 & 0.785 & 0.825 & 0.7847 & 0.8244 & +0.040 & +0.0397 \\
XGBoost & 0.8312 & 0.8662 & 0.790 & 0.800 & 0.7903 & 0.7972 & +0.010 & +0.0069 \\
Stacking & 0.8700 & 0.8738 & 0.855 & 0.840 & 0.8565 & 0.8410 & -0.015 & -0.0155 \\
\bottomrule
\end{tabular}
}
\end{table}

An ANOVA test across the 5-fold accuracy scores of all classifiers under the With-PCA setting gave
$F = 1.8205$, $p = 0.1052$, which is above $\alpha = 0.05$, so there is no statistically significant
difference between the classifiers once PCA is applied -- they all land in a fairly similar accuracy
band.

\section{Result Visualization}

\begin{figure}[H]
\centering
\includegraphics[width=\linewidth]{outputs/confusion_matrices_knn.png}
\caption{Confusion matrices for KNN, No-PCA (k=9, distance, euclidean) vs With-PCA (k=11, distance,
euclidean, 3 components), on the held-out test set.}
\end{figure}

\begin{figure}[H]
\centering
\includegraphics[width=\linewidth]{outputs/roc_curves_svm.png}
\caption{One-vs-rest ROC curves for the linear-kernel SVM, No-PCA (C=0.1) vs With-PCA (C=10, 3
components), on the held-out test set.}
\end{figure}

\textbf{Inference:} The With-PCA KNN model makes a couple fewer mix-ups between the two classes that
overlap the most, which lines up with it also scoring higher on test accuracy. The SVM ROC curves
hug the top-left corner in both versions with per-class AUC well above 0.9, so even where PCA costs a
few points of raw accuracy, the model's actual ability to tell the classes apart barely moves.

\section{Observations}
\begin{itemize}
\item \textbf{Which models improved most with PCA:} KNN, Gradient Boosting, Random Forest, AdaBoost
and XGBoost all picked up a bit of accuracy with the 3-component space, Gradient Boosting the most at
+0.04. Makes sense for the distance and tree-based methods -- once you strip out the correlated,
noisy dimensions, distance calculations get cleaner and trees don't have to work as hard to find good
splits.
\item \textbf{Which did not improve:} Naive Bayes took the biggest hit, down 0.05 in accuracy. That
tracks, since it assumes the features are independent, and PCA components (even though they're
decorrelated) reshuffle the variance in a way that doesn't play nicely with its per-feature Gaussian
assumption once you're down to just 3 of them. SVM, Logistic Regression, Decision Tree and Stacking
also dropped a little, somewhere between half a point and two and a half points.
\item \textbf{Did PCA make folds more stable:} Not consistently. Naive Bayes, Logistic Regression and
XGBoost actually got \emph{less} stable across folds after PCA, while SVM, AdaBoost and Random Forest
got more stable. So there isn't one clean story here, it really depends on the model.
\item \textbf{Did PCA help with overfitting on high-dimensional data:} With only 10 features to begin
with, this isn't really a high-dimensional problem, so I wouldn't expect a dramatic overfitting fix
either way. The real payoff from PCA here is more practical -- 10 features down to 3 while keeping
about 87\% of the variance, which is a nice speedup without hurting most models much.
\item \textbf{Linear models vs ensembles:} SVM and Logistic Regression both lost a little test
accuracy with PCA even though their CV means went up slightly, which suggests the reduced space
generalized fine on average across folds but happened to land a bit worse on this particular test
split. Ensembles like Random Forest, Gradient Boosting and XGBoost were more consistently the ones
that benefited, probably because fewer, decorrelated inputs make it easier for individual trees to
find clean splits.
\item \textbf{Was Stacking more robust:} Not really, at least not here. It had the best (or close to
best) No-PCA test accuracy out of all ten models at 0.855, but it also dropped the most among the
top performers once PCA was applied, down 0.015. It's basically just averaging out whatever its SVM,
Random Forest and XGBoost base learners were already doing, rather than being immune to PCA's effects
on its own.
\end{itemize}

\section{Discussion}
Putting all ten models side by side, PCA down to 3 components (about 87\% of the original variance)
tends to nudge cross-validation accuracy up a bit but nudge held-out test accuracy down for roughly
half the models, Naive Bayes and SVM being the clearest examples. The ANOVA result backs this up too
-- once PCA is applied, none of the ten classifiers is statistically distinguishable from any other on
accuracy, so at that point picking a model would come down to things like training time or how easy it
is to explain, not raw accuracy.

\section{Conclusion}
Honestly, PCA on this dataset comes out roughly accuracy-neutral overall. It clearly helps the
distance-based and some ensemble models (KNN, Random Forest, Gradient Boosting) by clearing out
redundant, correlated dimensions, but it quietly hurts models that either lean on feature independence
(Naive Bayes) or that were already doing fine finding a decent boundary in the full 10-dimensional
space (SVM, Logistic Regression, Decision Tree). The clearest win it gives you is the dimensionality
reduction itself, 10 features down to 3, and whether that's worth the small accuracy trade-off really
depends on which model you're using it with.

\section{Learning Outcomes}
\begin{itemize}
\item Got practice setting up a proper No-PCA vs With-PCA comparison with \texttt{Pipeline},
\texttt{GridSearchCV} and 5-fold CV across ten pretty different classifier families.
\item Learned to actually read a scree plot and decide on a component count by weighing variance
retained against what the models liked, instead of just chasing the usual 95\% rule blindly.
\item Saw directly that PCA isn't some universal upgrade -- whether it helps depends a lot on whether
the model cares about feature independence, distance, or a linear decision boundary.
\item Used ANOVA properly for once, to check whether a bunch of classifiers' CV scores are actually
different from each other rather than just eyeballing a table and guessing.
\item Got more comfortable stacking different kinds of base learners (SVM, Random Forest, XGBoost)
under a Logistic Regression meta-learner.
\end{itemize}

\section{References}
\begin{enumerate}
\item Stephen Marsland, ``Machine Learning - An Algorithmic Perspective'', 2nd Edition, Chapman and
Hall/CRC Machine Learning and Pattern Recognition Series, 2015.
\item Ethem Alpaydin, ``Introduction to Machine Learning'', 3rd Edition, The MIT Press, 2014.
\item Scikit-learn documentation: \url{https://scikit-learn.org/stable/modules/decomposition.html\#pca}
\end{enumerate}

\end{document}
